%%%%%%%%%%%%%%%%%%%%%%%%%%%%%%%%%%%%%%%%%%%%%%%%%%%%%%%%%%
% FILE: chapter04_floats.tex
%
% PURPOSE
% -------------------------------------------------------
% Explains how to insert figures, tables and algorithms
% using the high-level macros provided by the template.
%
%
% ARCHITECTURE ROLE
% -------------------------------------------------------
% User guide chapter.
%
% Documents the float wrapper macros defined in
% environments.tex.
%
%
% DEPENDENCIES
% -------------------------------------------------------
% environments.tex
% graphicx
% algorithm
%
%
% NOTES FOR DEVELOPERS
% -------------------------------------------------------
% If the float wrapper macros change, this chapter
% must be updated accordingly.
%
%%%%%%%%%%%%%%%%%%%%%%%%%%%%%%%%%%%%%%%%%%%%%%%%%%%%%%%%%%


\chapter{Floats de alto nivel}

La plantilla proporciona macros que simplifican la
inserción de figuras, tablas y algoritmos.

Estas macros encapsulan los entornos estándar de
\LaTeX{} y garantizan un formato consistente en
todo el documento.



% --------------------------------------------------
% CONCEPTO DE FLOAT
% --------------------------------------------------

\section{Concepto de float}

\begin{definitionbox}[Float]
En \LaTeX{}, un \textit{float} es un elemento del
documento que puede moverse automáticamente para
optimizar el diseño de la página.
\end{definitionbox}

Los floats más comunes son:

\begin{itemize}

\item figuras
\item tablas
\item algoritmos

\end{itemize}



% --------------------------------------------------
% INSERCIÓN DE FIGURAS
% --------------------------------------------------

\section{Inserción de figuras}

Las figuras se insertan mediante la macro:

\begin{lstlisting}[caption={Macro de inserción de figuras},label={code:figura}]
\figura{ruta_imagen}{caption}{fuente}{label}
\end{lstlisting}

\source{Plantilla}

Donde:

\begin{itemize}

\item \textbf{ruta\_imagen} es la ruta del archivo de imagen
\item \textbf{caption} es el título de la figura
\item \textbf{fuente} indica el origen de la figura
\item \textbf{label} es la etiqueta para referencias

\end{itemize}

\figura{03_figures/borrador.png}{Ejemplo de figura dentro del documento}{google imagenes}{fig:demo}



\begin{lstlisting}[caption={Ejemplo de inserción de figura},label={code:figuraej}]
\figura{03_figures/borrador.png}{Ejemplo de figura dentro del documento}{google imagenes}{fig:demo}
\end{lstlisting}

\source{Plantilla}



\begin{notebox}
Las imágenes utilizadas en el documento deben
almacenarse en la carpeta \texttt{03\_figures}.
\end{notebox}



% --------------------------------------------------
% INSERCIÓN DE TABLAS
% --------------------------------------------------

\section{Inserción de tablas}

Las tablas se insertan mediante la macro:

\begin{lstlisting}[caption={Macro de inserción de tablas},label={code:tabla}]
\tabla{contenido_tabla}{caption}{fuente}{label}
\end{lstlisting}

\source{Plantilla}

\tabla{\centering\begin{tabular}{cc}
\toprule
Dato 1 & Dato 2 \\
\midrule
Dato 3 & Dato 4 \\
\bottomrule
\end{tabular}}{Ejemplo de tabla}{Plantilla}{table:tablaej}

Ejemplo:

\begin{lstlisting}[caption={Ejemplo de código tabla},label={code:tablaej}]
\tabla{\centering\begin{tabular}{cc}
\toprule
Dato 1 & Dato 2 \\
\midrule
Dato 3 & Dato 4 \\
\bottomrule
\end{tabular}}{Ejemplo de tabla}{Plantilla}{tablaejemplo}
\end{lstlisting}

\source{Plantilla}



% --------------------------------------------------
% INSERCIÓN DE ALGORITMOS
% --------------------------------------------------

\section{Inserción de algoritmos}

Los algoritmos se insertan mediante la macro:

\begin{lstlisting}[caption={Macro de algoritmos},label={code:algoritmo}]
\algoritmo{contenido}{caption}{fuente}{label}
\end{lstlisting}

\source{Plantilla}



Ejemplo:

\begin{lstlisting}[caption={Ejemplo de algoritmo},label={code:algoritmoej}]
\algoritmo{
\begin{algorithmic}
\For{i = 1 to n}
\If{A[i] == x}
\Return i
\EndIf
\EndFor
\end{algorithmic}
}{Busqueda lineal}{Plantilla}{busqueda}
\end{lstlisting}

\source{Plantilla}

\algoritmo{
\begin{algorithmic}
\For{i = 1 to n}
\If{A[i] == x}
\Return i
\EndIf
\EndFor
\end{algorithmic}
}{Busqueda lineal}{Plantilla}{busqueda}

% --------------------------------------------------
% FUENTE DE FIGURAS
% --------------------------------------------------

\section{Fuente de figuras y tablas}

Todas las figuras y tablas incluyen automáticamente
una fuente mediante el comando:

\begin{lstlisting}[caption={Macro de fuente},label={code:source}]
\source{Plantilla}
\end{lstlisting}

\source{Plantilla}

Este comando genera una línea tipográfica centrada
debajo del caption que indica el origen del contenido.



% --------------------------------------------------
% VENTAJAS DEL SISTEMA
% --------------------------------------------------

\section{Ventajas del sistema}

Utilizar estas macros ofrece varias ventajas.

\begin{itemize}

\item formato uniforme en todo el documento
\item numeración automática
\item inclusión automática de fuentes
\item etiquetas consistentes

\end{itemize}



\begin{warningbox}
Insertar figuras o tablas directamente mediante
los entornos estándar de \LaTeX{} puede generar
inconsistencias de formato dentro del documento.
\end{warningbox}



\begin{examplebox}
Gracias a estas macros, el autor no necesita recordar
la estructura completa de los entornos \texttt{figure}
o \texttt{table}, lo que simplifica la escritura del
documento.
\end{examplebox}
